% SPDX-License-Identifier: CC-BY-4.0
% Copyright 2025 Florian Aram Feuerriegel — kassensturz.org
% Teil I — Das Modell in Worten

\section{Was für ein Modell das ist}

Die Engine ist ein \emph{Mikrosimulationsmodell mit makroökonomischem Rahmen}. Sie bildet
Deutschland durch zwölf typische Haushalte ab, je einen für die Einkommensdezile D1 bis D9 und
drei für das oberste Dezil (D10a: 90.--95.~Perzentil, D10b: 95.--99., D10c: das oberste
Prozent). Zusammen stehen sie für rund 41~Millionen Haushalte. Für jeden dieser Haushalte
rechnet die Engine Steuern, Beiträge, Transfers und Nettoeinkommen aus; die Summe über alle
Haushalte, gewichtet mit ihrer Zahl, ergibt die Einnahmen und Ausgaben des Staates.

Darüber liegt ein einfacher Makrorahmen: ein nominales Bruttoinlandsprodukt, das mit einem
Trend wächst, ein Schuldenstand, der sich aus dem Saldo fortschreibt, und ein
Emissionspfad. Die Zeit läuft in \emph{Perioden}, im Standardkurs fünf zu je vier Jahren
(2025--2044), entsprechend einer Legislaturperiode. Innerhalb einer Periode gilt eine Politik
unverändert; zwischen den Perioden wird der Zustand fortgeschrieben.

Drei Eigenschaften bestimmen, wie die Ergebnisse zu lesen sind:

\begin{enumerate}
\item \textbf{Alles ist ein Vergleich.} Die Engine sagt nicht, wie hoch das Nettoeinkommen
  eines Haushalts 2030 sein \emph{wird}, sondern wie es sich gegenüber dem Status quo
  \emph{derselben} Periode unterscheidet. Die Absolutwerte sind kalibriert, aber die Aussagen
  des Planspiels liegen in den Differenzen.
\item \textbf{Verhalten ist berücksichtigt, aber einfach.} Haushalte arbeiten etwas weniger,
  wenn ihnen von einem zusätzlichen Euro weniger bleibt; sie konsumieren weniger, wenn die
  Umsatzsteuer steigt; Unternehmen investieren weniger bei höheren Gewinnsteuern; Emissionen
  sinken mit dem \COz-Preis. Jede dieser Reaktionen ist eine einzige Elastizität aus der
  Literatur, keine geschätzte Nutzenfunktion. Es gibt keine Erwartungen, keine Optimierung
  und keine Rückwirkung auf Löhne oder Preise.
\item \textbf{Jede Zahl hat eine Quelle.} Elastizitäten, Faktoren und Tarifformeln stehen in
  den Objekten \datei{FORMEL\_QUELLEN\_*} und \datei{ELAST\_QUELLEN} mit Fundstelle. Der Anhang
  druckt sie vollständig ab. Wo keine Primärquelle erreichbar war, steht \enquote{Näherung}.
\end{enumerate}

\begin{merksatz}
Zwölf Haushaltstypen, ein Staat, ein Makrorahmen, Perioden zu vier Jahren. Ergebnisse sind
Unterschiede zum Status quo derselben Periode.
\end{merksatz}

\section{Eine Periode in zehn Schritten}
\label{sec:zehn}

Was geschieht, wenn ein Team eine Politik beschließt? Die Funktion \datei{berechne()} geht
für jede Periode dieselben zehn Schritte, Abbildung~\ref{fig:ablauf} zeigt sie im Überblick.

\begin{enumerate}[label=\arabic*.,leftmargin=2em]
\item \textbf{Geltendes Recht einrechnen.} Zu den Reglern kommen die beschlossenen Pfade
  hinzu, gemittelt über die Jahre der Periode: die Senkung der Körperschaftsteuer um je einen
  Punkt ab 2028 und der steigende Rentenbeitrag nach §\,158 SGB~VI. Die Regler zeigen die Werte
  von 2026.
\item \textbf{Arbeitsangebot.} Für jeden Haushalt wird der Grenzsteuersatz mit dem des Status
  quo verglichen. Bleibt vom letzten Euro mehr übrig, steigt das Arbeitseinkommen etwas.
\item \textbf{Einkommensteuer.} Auf das angepasste Arbeitseinkommen wird der Tarif des
  §\,32a EStG angewandt, auf Kapitaleinkommen die Abgeltungsteuer.
\item \textbf{Weitere Steuern.} Körperschaft- und Gewerbesteuer auf die Unternehmensgewinne,
  Umsatzsteuer auf den Konsum, Erbschaft-, Vermögen- und Bodensteuer, \COz-Abgabe auf die
  bepreisten Emissionen.
\item \textbf{Sozialbeiträge} auf die Löhne bis zur Beitragsbemessungsgrenze.
\item \textbf{Transfers und Renten:} Bürgergeld, Kindergeld, gegebenenfalls Klimageld; die
  Renten folgen dem Rentenniveau.
\item \textbf{Haushalte.} Aus allem entsteht das Nettoeinkommen jedes Haushaltstyps, dazu
  die Last aus Unternehmens- und Vermögensteuern, die am Ende Menschen tragen. Daraus folgen
  Gini, Palma und Armutsrisikoquote.
\item \textbf{Nachfrage.} Haushalte, die mehr (weniger) Geld haben als im Status quo, geben
  einen Teil davon aus (weniger aus). Zusammen mit zusätzlichen öffentlichen Investitionen
  verschiebt das die Wirtschaftsleistung dieser Periode.
\item \textbf{Fortschreibung in Euro der Periode.} Alles bis hierher ist in Größen von 2025
  gerechnet. Die Einnahmen werden mit dem tatsächlichen BIP der Periode hochgerechnet, die
  Ausgaben mit dem Preis- und Lohntrend. Dazu kommen Zinsen, Verteidigung und
  Sondervermögen.
\item \textbf{Saldo und Schuldenbremse.} Einnahmen minus Ausgaben ergeben den Saldo, bereinigt
  um Konjunktur und Ausnahmen den strukturellen Saldo, der gegen die Schuldenbremse geprüft
  wird.
\end{enumerate}

Danach übergibt \datei{berechneTransition()} an die nächste Periode: Das BIP wächst mit dem
Trend und den Angebotseffekten, die Schulden wachsen mit Zins und Defizit, die Emissionen
summieren sich zum verbrauchten \COz-Budget.

\begin{figure}[t]
\centering
\begin{tikzpicture}[
  node distance=3.5mm and 5mm,
  box/.style={draw=grau, rounded corners=2pt, fill=hell, align=center, font=\small,
              minimum height=8mm, text width=27mm},
  pfeil/.style={-{Stealth[length=2mm]}, thick, grau}]
\node[box] (a) {1 Geltendes\\Recht};
\node[box, right=of a] (b) {2 Arbeits-\\angebot};
\node[box, right=of b] (c) {3 Einkommen-\\steuer};
\node[box, right=of c] (d) {4 Weitere\\Steuern};
\node[box, below=of d] (e) {5 Sozial-\\beiträge};
\node[box, left=of e] (f) {6 Transfers\\und Renten};
\node[box, left=of f] (g) {7 Haushalte:\\Netto, Gini};
\node[box, left=of g] (h) {8 Nachfrage\\der Periode};
\node[box, below=of h] (i) {9 Fortschreibung\\nominal};
\node[box, right=of i] (j) {10 Saldo,\\Schuldenbremse};
\node[box, right=of j, fill=white, text width=58mm] (k) {Übergang zur nächsten Periode:\\BIP, Schulden, Emissionen, Kapital};
\draw[pfeil] (a)--(b); \draw[pfeil] (b)--(c); \draw[pfeil] (c)--(d); \draw[pfeil] (d)--(e);
\draw[pfeil] (e)--(f); \draw[pfeil] (f)--(g); \draw[pfeil] (g)--(h); \draw[pfeil] (h)--(i);
\draw[pfeil] (i)--(j); \draw[pfeil] (j)--(k);
\end{tikzpicture}
\caption{Ablauf einer Periode in \datei{berechne()} und Übergang in \datei{berechneTransition()}.}
\label{fig:ablauf}
\end{figure}

\section{Die Haushalte}

Jeder der zwölf Haushaltstypen ist durch wenige Zahlen beschrieben (Tabelle~\ref{tab:dezile}
im Anhang): Bruttoeinkommen, Anteil der Kapitaleinkommen, Konsumquote, Vermögen, Zahl der
Haushalte, das Bedarfsgewicht nach der neuen OECD-Skala und der Anteil der gesetzlichen Rente
am Einkommen. Grundlage sind SOEP v40, der Mikrozensus 2024, die Verteilungsrechnung DINA-DE
und der Vermögensbericht des DIW.

Das oberste Dezil ist dreigeteilt, weil sich dort fast alles abspielt, was Vermögen- und
Spitzensteuern betrifft: D10c, das oberste Prozent, hat im Mittel 700.000\,\euro{}
Bruttoeinkommen, zu 45\,\% aus Kapital, und 7~Mio.\,\euro{} Vermögen.

Zu jedem Haushaltstyp gibt es außerdem \emph{Profile}, die Haushalte und Staat gemeinsam
benutzen: Kinder mit Kindergeldanspruch (zusammen 17~Mio.), Bürgergeldbezug (zusammen
3,9~Mio. Regelsatz-Jahresäquivalente, nur D1 bis D4), die Last aus dem \COz-Preis
(regressiv: 4\,\% des Einkommens in D1, 1\,\% in D10c) und die Grenzkonsumneigung. Dass
beide Seiten dieselben Zahlen benutzen, ist nicht selbstverständlich: Bis September 2026
zählten die Haushalte 36,5~Mio. Kinder, der Staat 17~Mio. Seitdem prüft
\datei{tests/buchung.test.js}, dass die Summe über die Haushalte der Buchung des Staates
entspricht.

\begin{merksatz}
Zwölf Haushaltstypen mit Einkommen, Vermögen, Kindern, Transferbezug, Rente und
Konsumverhalten. Staat und Haushalte rechnen mit denselben Profilen.
\end{merksatz}

\section{Die Einkommensteuer}

Der Tarif ist der amtliche Formeltarif des §\,32a EStG 2026: steuerfrei bis zum
Grundfreibetrag (12.348\,\euro{}), dann eine Progressionszone, in der der Grenzsteuersatz
linear von 14 auf knapp 24\,\% steigt, eine zweite bis 42\,\%, eine Proportionalzone mit 42\,\%
und ab 277.826\,\euro{} der Spitzensatz von 45\,\%.

Vier Regler greifen in den Tarif ein, und jeder wirkt nur dort, wo sein Name hinzeigt: Der
\emph{Grundfreibetrag} verschiebt den ganzen Tarif, der \emph{Eingangssatz} hebt oder senkt den
Beginn der Progression, der \emph{Spitzensatz} gilt nur oberhalb der \emph{Grenze}, und die
Grenze verschiebt nur den Beginn der obersten Zone. Das klingt selbstverständlich, war es
aber nicht: Bis September 2026 hob ein höherer Spitzensatz die 42-\%-Zone mit an, und ein
Haushalt aus D5 zahlte für eine Reform der Spitzenverdiener mit.

Weil die Daten \emph{Haushalts}einkommen sind, der Tarif aber je Steuerpflichtigem gilt,
rechnet die Engine um: Nur 79\,\% des Bruttos sind zu versteuerndes Einkommen (Werbungskosten,
Vorsorgeaufwand), und ein Haushalt besteht im Mittel aus 1,6 Tarifeinheiten (Einzelpersonen
und zusammenveranlagte Paare). So trifft das Aufkommen die amtliche Größenordnung.

Für das oberste Prozent reicht ein Durchschnitt nicht: Sein mittleres zu versteuerndes
Einkommen liegt unter der Grenze des Spitzensatzes. Rechnete man nur mit dem Durchschnitt,
zahlte niemand im Modell den Spitzensatz. Die Engine nimmt deshalb an, dass die Einkommen
innerhalb von D10c einer Pareto-Verteilung folgen (Exponent 1,5, der schwere Rand der
Schätzungen für Deutschland), und berechnet, welcher Teil davon über der Grenze liegt.

\begin{merksatz}
Tarif 2026 nach §\,32a EStG, umgerechnet auf Haushalte. Jeder Regler wirkt nur auf seine
Zone. Das oberste Prozent ist als Verteilung modelliert, nicht als Durchschnitt.
\end{merksatz}

\section{Verhalten}

Vier Reaktionen machen aus der mechanischen Rechnung eine Verhaltensrechnung.

\textbf{Arbeitsangebot.} Steigt der Anteil, der vom letzten verdienten Euro übrig bleibt
(eins minus Grenzsteuersatz), um 10\,\%, steigt das Arbeitseinkommen um 2\,\%
(Elastizität 0,20, Konsens von Saez, Chetty und Gruber). Für das oberste Prozent ist die
Elastizität doppelt so hoch (0,40), weil dort Einkommen leichter verlagert werden kann;
oberhalb eines Grenzsteuersatzes von 45\,\% kommt Steuervermeidung hinzu, oberhalb von 60\,\%
Wegzug. Alle Reaktionen sind nach unten und oben begrenzt.

\textbf{Konsum.} Eine um einen Punkt höhere Umsatzsteuer senkt den besteuerten Konsum um
0,35\,\%, getrennt für den Regel- und den ermäßigten Satz.

\textbf{Investitionen.} Liegt die Summe aus Körperschaft- und Gewerbesteuersatz um zehn Punkte
über 30\,\%, sinken Gewinne und Investitionen um 4\,\% (Elastizität $-0{,}40$, Meta-Analyse
von Gechert und Heimberger). Das mindert die Bemessungsgrundlage sofort und das private
Kapital langfristig.

\textbf{Emissionen und Vermögen.} Ein um 100\,\euro{}/t höherer \COz-Preis senkt die bepreisten
Emissionen um 30\,\%. Eine Vermögensteuer senkt das angegebene Vermögen, nach Erfahrungen aus
der Schweiz (Brülhart u.\,a. 2022), hier halbiert.

\begin{merksatz}
Arbeit, Konsum, Investitionen, Emissionen und Vermögen reagieren über je eine Elastizität aus
der Literatur. Alle Reaktionen sind begrenzt, damit extreme Regler nicht ins Unsinnige
führen.
\end{merksatz}

\section{Die übrigen Einnahmen}

\textbf{Umsatzsteuer.} 70\,\% des Konsums fallen unter den Regelsatz, 30\,\% unter den
ermäßigten. Weil der Konsum der Haushalte nur einen Teil der Bemessungsgrundlage erfasst
(dazu kommen Käufe von Staat, Wohnungsbau und befreiten Branchen), wird hochgerechnet, und
3,7\,\% gehen durch Hinterziehung verloren (VAT Gap der EU-Kommission).

\textbf{Unternehmenssteuern.} Körperschaftsteuer (15\,\%) und Gewerbesteuer (effektiv 14\,\%)
haben getrennte Bemessungsgrundlagen, weil Personengesellschaften Gewerbe-, aber keine
Körperschaftsteuer zahlen. Beide schrumpfen mit der Investitionsreaktion.

\textbf{Vermögensbezogene Steuern.} Die Erbschaftsteuer rechnet mit 400\,\Mrd{} Erbmasse im
Jahr, von der nach den persönlichen Freibeträgen 45\,\% steuerpflichtig bleiben;
Betriebsvermögen ist weitgehend verschont. Eine Vermögensteuer trifft Vermögen über
2~Mio.\,\euro{} (3.500\,\Mrd), eine Bodenwertsteuer den Bodenwert (5.000\,\Mrd). Die
Mindeststeuer auf Milliardäre nach Zucman rechnet die gezahlte Einkommensteuer (rund 0,3\,\%
des Vermögens) und eine gleichzeitige Vermögensteuer an.

\textbf{\COz-Preis.} Die Emissionen folgen einem Basispfad nach dem Projektionsbericht 2025
des Umweltbundesamts: 649~Mt im Jahr 2025, $-63\,\%$ bis 2030 und $-80\,\%$ bis 2040 gegenüber
1990. Nur rund die Hälfte davon (327~Mt) trägt den nationalen \COz-Preis von 55\,\euro{}/t;
nur auf diesen Teil wirkt der Regler. Das Aufkommen fließt in den Klimafonds; ein Klimageld
gibt es im Status quo nicht, als Schalter aber schon: Es zahlt 70\,\% des Aufkommens pro
Haushalt gleich aus.

\textbf{Sozialbeiträge.} Renten-, Kranken-, Arbeitslosen- und Pflegeversicherung auf die
sozialversicherungspflichtige Lohnsumme von 1.750\,\Mrd, bis zur Beitragsbemessungsgrenze.
Der Beitragssatz wirkt \emph{nur auf die Einnahmen}. Das ist eine bewusste Entscheidung: Bis
September 2026 folgten auch die Ausgaben dem Satz, und eine Senkung des Rentenbeitrags
verbesserte den Saldo und machte alle Haushalte reicher.

\begin{merksatz}
Jede Steuerart hat ihre eigene Bemessungsgrundlage, kalibriert auf das Ist-Aufkommen.
Beitragssätze bestimmen Einnahmen, nicht Leistungen.
\end{merksatz}

\section{Ausgaben, Transfers und Rente}

Die Ausgaben des Staates beginnen mit dem Stand von 2025, grob nach Aufgaben gegliedert
(soziale Sicherung, Gesundheit, Bildung, Verteidigung, Infrastruktur, Verwaltung, Zinsen).
Die Stellgrößen ändern daran:

\begin{itemize}
\item \textbf{Bürgergeld und Kindergeld:} Regelsatz beziehungsweise Betrag je Kind, mal
  Zahl der Berechtigten aus den Haushaltsprofilen, mal zwölf Monate. Unterkunft und
  Mehrbedarfe (10,8\,\Mrd) bleiben fest.
\item \textbf{Rentenniveau:} Die Rentenzahlungen (rund 362\,\Mrd) steigen oder fallen mit dem
  Sicherungsniveau (Status quo 48\,\%). Jede Änderung kommt in genau der Höhe bei den
  Haushalten an, die der Staat bucht, verteilt nach dem Rentenanteil ihres Einkommens.
  Die Alterung erhöht die Rentenausgaben von allein, bis 2041 um ein Viertel.
\item \textbf{Klimageld}, falls eingeschaltet.
\item \textbf{Erhebungskosten:} Ändert sich das Steuersystem, ändern sich die Kosten, es zu
  erheben (etwa 2,5\,\% des Einkommensteueraufkommens, 8\,\% der Erbschaftsteuer). In den
  Saldo geht nur die Änderung gegenüber dem Status quo.
\end{itemize}

Dazu kommt geltendes Recht, das über die Jahre wirkt: Die Verteidigungsausgaben steigen nach
dem Plan der Bundesregierung von 2,4 auf 3,5\,\% des BIP (2029), das Sondervermögen
Infrastruktur gibt von 2026 bis 2036 jährlich 45\,\Mrd{} aus, und der Klimafonds schrumpft
mit den Emissionen, aus deren Bepreisung er sich speist.

\begin{merksatz}
Ausgaben folgen Leistungen, nicht Beiträgen. Was der Staat für Haushalte ausgibt, kommt dort
in derselben Höhe an.
\end{merksatz}

\section{Vom Staat zum Haushalt: wer trägt was}

Das Nettoeinkommen eines Haushalts ist
\begin{center}
Brutto $-$ Einkommensteuer $-$ Beiträge $-$ Umsatzsteuer $-$ \COz-Last $+$ Transfers $+$
Rentenänderung $-$ Last aus Unternehmens- und Vermögensteuern.
\end{center}

Zwei Posten verdienen eine Erklärung. \emph{Wer trägt eine Beitragserhöhung?} Formal zahlen
Arbeitnehmer und Arbeitgeber je die Hälfte. Langfristig wälzen Arbeitgeber ihren Anteil
aber über niedrigere Löhne auf die Beschäftigten ab (Meta-Analyse von Melguizo und
González-Páramo). Die Engine lässt deshalb eine \emph{Änderung} der Beitragssätze voll bei
den Haushalten ankommen.

\emph{Wer trägt die Körperschaftsteuer?} Unternehmen sind keine Menschen; ihre Steuern tragen
am Ende Beschäftigte über niedrigere Löhne oder Eigentümer über niedrigere Gewinne. Nach
Fuest, Peichl und Siegloch (2018) tragen die Beschäftigten rund die Hälfte der
Gewerbesteuer. Die Engine verteilt eine Änderung von Körperschaft- und Gewerbesteuer je zur
Hälfte nach Arbeits- und Kapitaleinkommen, Erbschaft-, Vermögen- und Bodensteuer nach dem
Vermögen, die Milliardärssteuer auf D10c.

Aus den zwölf Nettoeinkommen entstehen die Verteilungsmaße. Wie in der amtlichen Statistik
(EU-SILC) werden sie auf das \emph{Äquivalenzeinkommen} bezogen: Ein Haushalt mit vier
Personen braucht mehr als einer mit einer, also wird durch das Bedarfsgewicht geteilt.
Der \emph{Gini-Koeffizient} misst die Fläche zwischen Lorenzkurve und Gleichverteilung (0 =
alle gleich, 1 = einer hat alles), die \emph{Palma-Ratio} teilt das Einkommen der oberen
10\,\% durch das der unteren 40\,\%. Die \emph{Armutsrisikoquote} zählt, wer weniger als 60\,\%
des mittleren Einkommens hat; weil ein Durchschnitt je Dezil dafür zu grob ist, wird der
Anteil innerhalb der unteren drei Dezile aus dem Status quo fortgeschrieben.

\begin{merksatz}
Jede Steuer und jede Leistung landet bei einem Haushaltstyp. Verteilungsmaße werden wie in
der amtlichen Statistik auf Äquivalenzeinkommen gerechnet.
\end{merksatz}

\section{Nachfrage}

Wenn Haushalte mehr Geld haben, geben sie einen Teil davon aus, und dieser Teil ist unten
größer als oben: Ein Haushalt in D1 gibt von einem zusätzlichen Euro 65~Cent aus, einer in
D10c 15~Cent (Grenzkonsumneigung nach Jappelli und Pistaferri). Die Engine summiert die
Änderung der Nettoeinkommen gegenüber dem Status quo, gewichtet mit dieser Neigung, und
rechnet sie mit dem Steuer- und Transfermultiplikator (0,6 bei mittlerer Neigung, Gechert
2015) in eine Nachfrageänderung um. Zusätzliche öffentliche Investitionen wirken mit dem
Multiplikator 1,0.

Die Nachfrageänderung wirkt \emph{nur in der laufenden Periode}: Sie hebt oder senkt das BIP
und über das BIP die Einnahmen (etwa 37~Cent je Euro, der automatische Stabilisator). Sie
wirkt symmetrisch: Wer konsolidiert, zahlt mit Wachstum; wer entlastet, gewinnt welches.
Dauerhaft bleibt davon nur, was in öffentliches Kapital fließt.

\begin{merksatz}
Entlastung unten wirkt stärker auf die Nachfrage als oben. Nachfrage wirkt eine Periode,
Kapital bleibt.
\end{merksatz}

\section{Von Periode zu Periode}

\textbf{Wirtschaftsleistung.} Das nominale BIP wächst im Trend um 2,5\,\% im Jahr (2\,\%
Inflation, 0,5\,\% real). Dazu kommen drei Niveaueffekte der Angebotsseite:
\begin{itemize}
\item \emph{Arbeit}: Mehr Arbeitsangebot hebt das BIP sofort, mit dem Arbeitsanteil 0,65.
\item \emph{Privates Kapital}: Höhere Gewinnsteuern senken das Kapital langfristig (mit dem
  Kapitalanteil 0,35), aber langsam: Der Kapitalstock passt sich um 7\,\% im Jahr an.
\item \emph{Öffentliches Kapital}: Zusätzliche öffentliche Investitionen, auch das
  Sondervermögen, bauen einen Kapitalstock auf, der mit 4\,\% im Jahr abschreibt. Ein
  zusätzlicher Stock in Höhe des heutigen öffentlichen Kapitals (1.600\,\Mrd) hebt das BIP um
  8\,\% (Bom und Ligthart 2014).
\end{itemize}
Alle drei sind Niveaus, keine Wachstumsraten: Eine höhere Körperschaftsteuer senkt das BIP
auf ein niedrigeres Niveau, aber nicht jede Periode weiter.

\textbf{Schulden.} Die Schuldenstandsquote wächst jedes Jahr mit dem Zins von 2,0\,\% und sinkt
mit jedem Primärüberschuss; das BIP im Nenner wächst mit 2,5\,\%. Weil der Zins unter der
Wachstumsrate liegt, stabilisiert sich eine Schuldenquote schon bei einem kleinen
Primärdefizit. Ob \enquote{Sparen alternativlos} ist, entscheidet das Spiel also nicht vorab.

\textbf{Emissionen.} Die jährlichen Emissionen folgen dem Basispfad, verringert um die
Wirkung der Politik, und summieren sich Jahr für Jahr. Dem steht ein deutsches Restbudget von
5.380~Mt gegenüber (Bevölkerungsanteil am globalen Budget für 1,7\,°C nach Forster u.\,a.
2026). Einen Klimaschaden auf das deutsche BIP rechnet das Modell nicht: Deutschland
verursacht 1,5\,\% der Weltemissionen, der Rückkanal auf die eigene Wirtschaft ist
vernachlässigbar. Das Klimaziel steht über Emissionen und Budget im Spiel.

\textbf{Demografie und Schocks.} Ein Rentenlastfaktor nach der 14. Bevölkerungsvorausberechnung
erhöht die Rentenausgaben (1,00 in 2025, 1,25 in 2041). Die Lehrperson kann Schocks setzen:
einen BIP-Einbruch, einen Sprung der Schuldenquote oder einen Zinsanstieg.

\begin{merksatz}
Das BIP wächst mit Trend, Arbeit, privatem und öffentlichem Kapital, alles als Niveaus. Die
Schulden wachsen mit Zins und Defizit. Die Emissionen summieren sich gegen ein Budget.
\end{merksatz}

\section{Schuldenbremse und Tragfähigkeit}

Seit der Reform von 2025 erlaubt die Schuldenbremse Bund und Ländern zusammen ein
strukturelles Defizit von 0,70\,\% des BIP. Strukturell heißt: ohne Konjunktureinfluss (ein
Prozent Nachfragelücke bewegt den Saldo um 0,5\,\% des BIP, Budget-Semielastizität der
EU-Kommission) und ohne die Ausnahmen (Verteidigung über 1\,\% des BIP, Sondervermögen). Im
Status quo der ersten Periode liegt der strukturelle Saldo bei \SQstruktur\,\% des BIP; die
Schuldenbremse ist damit verletzt. Das Modell rechnet den Gesamtstaat, nicht Bund und Länder
getrennt.

Dazu kommen Tragfähigkeitsindikatoren aus \datei{abgeleitet.js}: der Primärsaldo, der
Primärüberschuss, der die Schuldenquote gerade stabilisiert (Domar-Bedingung), die Lücke
zwischen beiden und ein Index der Generationengerechtigkeit aus Schulden und verbrauchtem
\COz-Budget.

\section{Der Status quo als Referenz}

Tabelle~\ref{tab:verlauf} zeigt, wohin der Status quo ohne jede Politikänderung führt. Er ist
nicht stabil: Demografie, Verteidigung und Zinsen treiben das Defizit, und die
Schuldenquote steigt von 63,5\,\% auf \SQschuldEnde\,\%. Die Teams entscheiden also nicht
zwischen \enquote{weiter so} und Veränderung, sondern darüber, wer die Anpassung trägt.

\begin{table}[!htb]
\centering\small
% Erzeugt von docs/modell/messung.js — nicht von Hand bearbeiten
\begin{tabular}{lrrrrrrrr}
\toprule
Periode & BIP & Schulden & Saldo & Saldo & strukt. & Emiss. & CO\textsubscript{2} kum. & Gini \\
 & Mrd.\,\euro{} & \% BIP & Mrd.\,\euro{} & \% BIP & \% BIP & Mt/a & Mt & \\
\midrule
2025-$-$1028 & 4.470 & 63,5 & $-$173,0 & $-$1,87 & $-$1,28 & 649 & 0 & 0,312 \\
2029-$-$1032 & 4.968 & 71,5 & $-$140,7 & $-$1,85 & $-$1,43 & 500 & 2.373 & 0,312 \\
2033-$-$1036 & 5.529 & 82,2 & $-$189,4 & $-$1,23 & $-$1,91 & 399 & 4.199 & 0,312 \\
2037-$-$1040 & 6.143 & 93,4 & $-$191,5 & $-$1,75 & $-$1,25 & 314 & 5.669 & 0,312 \\
2041-$-$1044 & 6.771 & 102,5 & $-$165,4 & $-$1,40 & $-$1,90 & 238 & 6.799 & 0,312 \\
\bottomrule
\end{tabular}

\caption{Status quo über fünf Perioden. BIP und Schuldenquote zu Beginn der Periode, übrige
Größen als Jahreswert der Periode, \COz{} kumuliert seit 2025 zu Periodenbeginn.}
\label{tab:verlauf}
\end{table}

\section{Was das Modell nicht kann}
\label{sec:grenzen}

Ein Lehrmodell muss vereinfachen. Diese Vereinfachungen sollte kennen, wer die Ergebnisse
erklärt:

\begin{itemize}
\item \textbf{Kein allgemeines Gleichgewicht.} Löhne, Preise und Zinsen reagieren nicht auf
  die Politik. Der Zins ist fest (außer im Zinsschock).
\item \textbf{Zwölf Haushaltstypen statt Millionen Haushalten.} Innerhalb eines Dezils sind
  alle gleich, außer im obersten Prozent (Pareto-Rand) und bei der Armutsquote.
\item \textbf{Näherungen} ohne Primärtabelle, weil die Quellen beim Anlegen nicht erreichbar
  waren: Rentenanteile je Dezil, Profil der Grenzkonsumneigung, öffentlicher Kapitalstock,
  Milliardärsvermögen, Pareto-Exponent. Sie sind im Code als solche markiert.
\item \textbf{Restabweichung Haushalt gegen Staat.} Bei Umsatzsteuer und Beiträgen weicht die
  Summe über die Haushalte noch um bis zu 9\,\% von der Buchung des Staates ab.
\item \textbf{Nicht gerechnete Schockwirkungen.} Zwei der fünf Schockeffekte
  (Investitionsrückgang, Emissionsminderung) stehen in der Bibliothek, wirken aber nicht; die
  Oberfläche kennzeichnet sie.
\item \textbf{Schuldenbremse} für den Gesamtstaat, nicht getrennt nach Bund und Ländern;
  EU-Fiskalregeln fehlen.
\end{itemize}

Die Messung für Teil~III hat drei weitere Eigenschaften sichtbar gemacht:

\begin{itemize}
\item \textbf{Die Beitragsbemessungsgrenze wirkt asymmetrisch.} Eine Anhebung um 10.000\,\euro{}
  verbessert den Saldo der ersten Periode um rund 6\,\Mrd, eine Senkung um denselben Betrag
  \emph{verbessert} ihn ebenfalls, um rund 3,5\,\Mrd. Grund: Beim Staat sind Mehreinnahmen
  aus der Grenze nach unten bei null abgeschnitten, bei den Haushalten sinken die Beiträge.
  Die entlasteten Haushalte fragen mehr nach, und der Stabilisator hebt den Saldo. Das ist
  ein Buchungsfehler; er ist zur Behebung vorgemerkt.
\item \textbf{Der Spitzensteuersatz bringt wenig:} rund \wert{zerl}{spitze}{gesamt}\,\Mrd{} je
  Prozentpunkt in der ersten Periode, weniger als übliche Schätzungen (etwa 1 bis 2\,\Mrd).
  Die Zerlegung in Abschnitt~\ref{sec:zerlegung} zeigt, dass das fast ganz an der Rechnung
  selbst liegt (\wert{zerl}{spitze}{mechanisch}\,\Mrd{} mechanisch), nicht am Verhalten
  (\wert{zerl}{spitze}{verhalten}\,\Mrd): Der Satz trifft nur den Teil der Einkommen im
  obersten Prozent, der nach dem Pareto-Rand über der Grenze liegt. Ob der Rand zu dünn ist,
  wäre mit der Einkommensteuerstatistik zu prüfen.
\item \textbf{Die Mindeststeuer auf Milliardäre greift erst oberhalb von 0,3\,\%}, weil die
  gezahlte Einkommensteuer angerechnet wird. Das ist gewollt, erklärt aber, warum kleine
  Sätze nichts bewirken.
\end{itemize}
